\chapter{Executive Summary}
\label{ch:executive_summary}

The recruitment industry heavily relies on \textbf{Applicant Tracking Systems (ATS)} that use rigid keyword-matching paradigms. Unfortunately, this means highly qualified candidates are routinely rejected simply because their resumes do not mirror the exact vocabulary of a job description. To solve this, I designed and validated an automated, human-centric resume screening pipeline that replaces brittle keyword filters with \textbf{Natural Language Processing (NLP)} and \textbf{Learning-to-Rank (LTR)} machine learning.

My research tackles the challenge of candidate--job fit from two fundamental angles:
\begin{itemize}
    \item \textbf{Semantic Understanding:} I contrasted traditional sparse statistical arrays (TF-IDF) against \textbf{dense semantic tensors (Sentence-BERT)}. This allows the system to understand the actual meaning behind a candidate's skills, rather than just counting overlapping words.
    \item \textbf{Human Intuition Modeling:} Beyond text, I formalized how a candidate's ``Experience Gap'' (the difference between their years of experience and the job's requirement) should impact their score. I compared a linear baseline and a piecewise quadratic penalty against a novel \textbf{Sigmoid-Bayesian posterior confidence framework}. I engineered this Bayesian approach specifically to replicate human judgment in the absence of massive historical hiring datasets.
\end{itemize}

To guarantee that this system works in the real world, I eliminated target leakage---a critical risk where models just memorize historical job buzzwords. I evaluated every model using a strict \textbf{\texttt{GroupShuffleSplit}} architecture, ensuring that test-time job postings were completely unseen during the training phase. 

I then benchmarked three different modeling paradigms head-to-head: continuous regression, binary classification, and specialized Learning-to-Rank estimators. The results were unambiguous: regression and classification both fail to produce what a recruiter actually needs---a ranked shortlist.

\textbf{Key Research Achievements:}
\begin{itemize}
    \item \textbf{Optimal Architecture:} I identified \textbf{\texttt{LGBMRanker}}, driven by LambdaMART gradients over dense SBERT embeddings, as the superior framework.
    \item \textbf{Peak Accuracy:} The system achieved an outstanding Normalized Discounted Cumulative Gain in the top 10 results (\textbf{NDCG@10}) of \textbf{0.9398}, proving the model's top recommendations closely match ideal human-curated shortlists.
    \item \textbf{Baseline Superiority:} The listwise ranking approach structurally outperformed both continuous regression ($R^2 = 0.3202$) and binary classification (ROC-AUC $= 0.8121$) baselines.
\end{itemize}

These findings confirm that candidate screening is fundamentally a \textbf{listwise sorting problem}, not an absolute scoring or pass/fail task. By deploying gradient-boosted ranking algorithms over dense semantic embeddings, we can directly resolve the limitations of deterministic ATS platforms. 

Beyond the statistical wins, this pipeline delivers tangible operational value. In my study corpus, an average job posting attracted roughly 330 applications. By accurately narrowing this down to a high-precision \textbf{top-10 shortlist}, this system significantly reduces the manual screening burden on HR professionals. In Section~\ref{ch:recommendation} (Recommendation), I outline the concrete steps---including computational latency planning, counterfactual bias mitigation, and business-intelligence integration---required to successfully transition this research prototype into a \textbf{production-ready enterprise system}.